\documentclass[12pt,a4paper]{article}
\usepackage[margin=24mm]{geometry}
\usepackage[T1]{fontenc}
\usepackage{lmodern,amsmath,amssymb,booktabs,array,xcolor}
\usepackage[colorlinks=true,urlcolor=blue!55!black,linkcolor=blue!55!black,citecolor=blue!55!black]{hyperref}
\setlength{\parindent}{0pt}
\setlength{\parskip}{6pt}
\setlength{\emergencystretch}{2em}
\title{Why the Compiled 4D Prefix\protect\\Implementation Is Faster}
\author{Implementation audit and reproducible 4D benchmarks}
\date{4 October 2026}
\begin{document}
\maketitle
\begin{abstract}
The compiled prefix implementation in \texttt{HV4DMagnitude} is faster than
the new numerical Chan backend on the 4D inputs tested here. On spherical
fronts with 20 or 40 points it is 1.59--2.11 times faster than the existing
Python Chan $d/3$ implementation. At 20 points it is 84.59--96.99 times faster
than the new default Numba backend. Two controlled comparisons explain this:
the old implementation compiles the entire numerical $2+2$ loop, while the new
compression policy creates costly intermediate term states whose manipulation
largely remains in Python. This report separates empirical speed, correctness,
and worst-case complexity. No solver was modified, and only 4D was tested.
\end{abstract}

\section{Question and measured answer}
The earlier comparison of numerical Chan implementations omitted the compiled
prefix solver in \texttt{HV4DMagnitude}. It therefore did not assess the user's
earlier Numba approach. We now compare that solver~\cite{old} with the new
numerical engine and the two existing Python Chan implementations~\cite{chan,new}.
The user's recollection of a faster compiled 4D implementation is confirmed.

\begin{table}[ht]
\centering\small
\setlength{\tabcolsep}{4pt}
\begin{tabular}{lrrrrr}\toprule
Family seed / $n$ & Old JIT & Chan $d/3$ & New JIT & No-comp. & Chan $d/2$\\\midrule
Sphere 42 / 10 & 5.295 & 7.519 & 9.259 & 10.972 & 1.515\\
Sphere 42 / 20 & 21.521 & 35.635 & 2087.270 & 45.190 & 5.167\\
Sphere 42 / 40 & 62.343 & 111.593 & budget & 142.064 & 19.298\\
Sphere 99 / 20 & 21.924 & 34.805 & 1854.674 & 37.848 & 5.983\\
Sphere 99 / 40 & 52.861 & 111.451 & budget & 157.691 & 21.045\\
Grid 7 / 20 & 0.834 & 2.225 & 5.549 & 5.907 & 0.323\\
Uniform 17 / 20 & 14.546 & 18.634 & 832.086 & 29.710 & 2.405\\
\bottomrule\end{tabular}
\caption{Median of three warm end-to-end calls, in milliseconds. ``budget''
denotes a 35-second whole-worker timeout, not a per-call lower bound.}
\label{tab:main}
\end{table}

Old JIT uses \texttt{PrefixHV4D} with \texttt{base\_hard=2} and
\texttt{fast=True}.

New JIT uses \texttt{NumericalChanNumba} with
dimension four and the same base threshold, the shared backend used by
the new 4D wrapper. No-comp. is that same new backend with
\texttt{block\_levels=10**6}; this diagnostic avoids compression on these
inputs but has no claimed worst-case guarantee. Chan $d/3$ and $d/2$ use no Numba.

\textbf{Chan $d/2$ is fastest on all seven measured inputs.} On the four
spherical $n=20,40$ cases, it is 2.51--4.16 times faster than old JIT.
These are small input sizes: this observation neither contradicts an asymptotic
advantage for a different algorithm nor establishes an eventual crossover.
No exponent is fitted to this experiment.

\section{Fair inputs, timing, and validation}
All algorithms compute the same ordinary four-dimensional hypervolume.
Old prefix and both Chan references receive nonnegative minimization points
$p$ with reference $r$. The new engine receives the reflected anchored corners
$r-p$. Passing $-p$ with zero reference to the old prefix solver would be
unsafe: its grids are anchored at zero and discard negative cuts.
The conversion for the new engine is included in its timed call.
No implementation uses a dominance prefilter, and the shared base threshold is
two hard boxes wherever that option exists.

Each algorithm and dataset uses a fresh process. One initial call precedes
three timed calls; the table reports their median. Imports and the initial
call are excluded, but packing, preprocessing, and complete evaluation are
included. The initial call may load a Numba disk cache; its recorded duration
is not a pure compilation time. Workers run sequentially in a reproducibly
shuffled order within each dataset. There is no Numba parallelism or fastmath.
A timeout includes process startup, imports, first use and all three repeats.
It gives no measured per-call ratio for an unfinished row.

Sphere inputs are absolute independent Gaussian vectors normalized to unit
Euclidean length, with $r=(1.2,1.2,1.2,1.2)$. Seed 42 consumes sizes
$5,10,20,40$ successively, discarding the five-point case; seed 99 consumes
$20,40$. Grid inputs use independent integers from zero through three divided
by four; uniform inputs use independent coordinates in $[0,1)$. Both have
$r=(1,1,1,1)$. Every input is saved, rather than reconstructed from timings.
The historical audit did not archive its exact timing inputs: this experiment
uses its committed test family, without claiming to reproduce its historical
table byte for byte.

\textbf{All 52 completed rows agree numerically.} Against the separate Chan
$d/2$ evaluator, maximum scaled disagreement is $7.89\times10^{-16}$, using
$|x-y|/\max(1,|y|)$. The ten-point sphere also passes an independent
inclusion--exclusion oracle evaluated with exact rational representations of
the supplied binary floating-point inputs; its maximum scaled error is
$2.22\times10^{-16}$. Four new default runs time out: Python and Numba on
each of the two forty-point spheres. These checks are empirical validation,
not a general numerical-error theorem. Every completed old-JIT row records
a Numba nopython signature.

The machine is an AMD Ryzen 7 5800H with eight physical cores, sixteen logical
processors, and approximately 16 GiB RAM. Software is Windows 11,
CPython 3.12.14, NumPy 2.5.3 and Numba 0.68.0. The measurements do not use
confidence intervals, processor pinning, or a controlled dedicated server;
three repeats support the large observed differences, not fine percentage claims.

\section{The compilation boundary matters}
The old implementation packs six staircases, reversed values, target grids,
and cumulative source records into numeric arrays. Its
\texttt{fastkernel.k\_prefix} compiles the entire nested target-plane scan:
binary searches, clipping source bounds, choosing the source-area formula,
and numerical accumulation all execute inside the compiled call.

Disabling JIT \emph{for that same flat kernel} is the cleanest compilation
comparison. Table~\ref{tab:python} also includes the old object-based evaluator,
which is a distinct code path selected by \texttt{fast=False}.

\begin{table}[ht]
\centering\small
\setlength{\tabcolsep}{5pt}
\begin{tabular}{lrrrr}\toprule
Family seed / $n$ & Old flat Python & Old object Python & Flat/JIT & New Python\\\midrule
Sphere 42 / 10 & 69.091 & 69.840 & 13.05 & 8.843\\
Sphere 42 / 20 & 255.417 & 250.038 & 11.87 & 1291.028\\
Sphere 42 / 40 & 804.383 & 785.652 & 12.90 & budget\\
Sphere 99 / 20 & 339.409 & 330.800 & 15.48 & 1147.316\\
Sphere 99 / 40 & 695.438 & 715.985 & 13.16 & budget\\
Grid 7 / 20 & 1.596 & 2.260 & 1.91 & 3.424\\
Uniform 17 / 20 & 139.630 & 129.447 & 9.60 & 500.031\\
\bottomrule\end{tabular}
\caption{Warm medians in milliseconds, except the dimensionless Flat/JIT
ratio. ``budget'' has the same meaning as in Table~\ref{tab:main}.}
\label{tab:python}
\end{table}

On the spherical $n=20,40$ cases, compiling the old flat path gives an
11.87--15.48-fold speedup. The new backend instead compiles small helpers for
searches, prefix sums, normalization, masks and related array operations.
Winner enumeration and much of the copying, comparison and merging of terms
remain in Python. Numerous calls on short arrays do not remove that work.
``New Python'' in Table~\ref{tab:python} means the shared numerical engine
at dimension four, not the separate specialized \texttt{numerical\_chan4.py}
file. All these timing inputs use floating-point arithmetic, including this
otherwise exact-capable Python engine.

\section{Compression changes the later cost of evaluation}
The new default schedule compresses after a block of four cyclic levels at
these input sizes. In 4D this temporarily creates eight variables and eliminates
four, producing signed terms. The old prefix solver performs no such compression.
The existing Python Chan $d/3$ code additionally gates compression on representation
size; it performs none on the seed-42 twenty-point case.

For that case the two new Numba configurations have the same 61 recursive
calls and 30 cuts, but very different term work:

\begin{center}\small
\begin{tabular}{lrr}\toprule
Statistic & Default & No compression\\\midrule
Compressions & 10 & 0\\
Raw emitted terms & 10,809 & 709\\
Term eliminations & 7,441 & 526\\
Peak intermediate terms & 625 & 11\\
Median milliseconds & 2,087.270 & 45.190\\\bottomrule
\end{tabular}
\end{center}

The no-compression variant is 46.19 times faster on this case and 49.00 times
faster on the other twenty-point sphere. This measures the compression policy
\emph{and the states it leaves for later work}, not merely the cost of calling
the compression routine. The ten-point sphere and tied-grid case perform no
compression under either policy, so their small timing differences are not
evidence of a compression effect.

Separate profiling of the seed-42 twenty-point case supports this explanation.
The old JIT path records 40,833 Python-visible calls; the new default records
3,285,965, and the new no-compression variant 82,860. Calls internal to compiled
code are not counted individually, so these figures describe interpreter
overhead, not arithmetic-operation complexity. In the new default profile,
about 15.3\% of time is cumulative within \texttt{push\_blocks}, while 78.7\%
is in base-case evaluation. Expensive states persist beyond compression itself.
Prominent Python-side work includes \texttt{emitted}, \texttt{impose},
\texttt{add\_edge}, unary multiplication and merging. Profiled durations
are never substituted for the unprofiled benchmark medians.

The old solver also absorbs slabs by shrinking geometric cells and reclassifying
the remaining orthants; the new engine zeros unary entries without shrinking
the cell. This is a further difference, but its individual speed contribution
was not isolated here. Cross-implementation counters are not universally
comparable: the new node count includes axis-skip calls and its term high-water
mark includes temporary elimination outputs.

\section{What the numerical $2+2$ loop provides}
The six pairwise staircases encode constraints between coordinates. After
fixing a target cell $(u,v)$ in coordinates 3 and 4, cross constraints give
caps $A(u,v),B(u,v)$ on coordinates 1 and 2. The remaining source-plane area
$H_{12}(A,B)$ is obtained from the stored records for the first staircase.
The complete computation has the numerical form
\[
 K=\sum_{u,v} w_3(u)w_4(v)
       \mathbf{1}_{\{v<f_{34}(u)\}}
       H_{12}\bigl(A(u,v),B(u,v)\bigr).
\]
Here $u,v$ denote representatives of target cells and the $w_i$ are their
widths for ordinary volume. The source evaluator uses a two-branch prefix
formula; no symbolic term is constructed for each target cell. In the
seed-42 twenty-point run there are 81 box evaluations, each making sixteen
signed prefix queries, for 1,296 queries in total.

This is a useful compact numerical implementation, but it explicitly visits
target-grid cells at every prefix query. It also retains live orthants in
descendants and does not implement the new compression schedule.
\textbf{This audit does not establish an overall
$O(n^{4/3}\operatorname{polylog}n)$ bound for that older implementation.}
Its own repository marks the overall complexity claim as unverified.
The fast constant factors observed here do not supply the missing bound on
global cell visits. Conversely, disabling the new engine's compression is
only an ablation: its speed does not preserve a guarantee proved for a different
schedule.

\section{Magnitude: a separate correctness finding}
The old public hypervolume solver hardcodes Lebesgue measure, so the main
benchmark concerns ordinary HV. Direct calls to its packed kernel with the
magnitude product measure reveal a separate atom-handling defect.
Let all six staircases be constant two, with every axis grid equal to $[0,1]$.
The magnitude of the four-dimensional unit box should be
$(1+1/2)^4=81/16=5.0625$. The decomposed and brute evaluators return this value,
but the packed kernel returns $2.25$. At query $(1,1,0,0)$, expected magnitude
is $9/4$, while the packed kernel returns zero. Lebesgue versions of these
two checks return the correct values, one and zero.

The target-grid loop substitutes the origin atom for the first interval
instead of accounting for both pieces. This defect does not explain the
ordinary-HV speed difference. The ordinary-HV route already supports magnitude
through the affine identity, for nonnegative anchored generators,
\[
 \operatorname{Mag}(D(A))=
 \operatorname{HV}_4\!\left(\{(1+a_1/2,\ldots,1+a_4/2):a\in A\}\right).
\]
To see it, each nonempty intersection of anchored boxes has coordinatewise
minimum $b$ and magnitude $\prod_i(1+b_i/2)$. Its transformed ordinary volume
is the same product, since the affine coordinate transformation commutes with
minimum. Inclusion--exclusion proves equality for the union. This is an
equivalence of numerical set values; it is not a translation of the anchored
atomic measure.

Three exact inclusion--exclusion spot checks of this route pass through the
old public HV solver: the unit box gives $81/16$, the degenerate box with
corner $(1,1,0,0)$ gives $9/4$, and the four-generator set
\[
 \{(1,4,3,2),(2,3,4,1),(3,2,1,4),(4,1,2,3)\}
\]
gives $765/16$. The direct packed-kernel defect is recorded in the accompanying
reproducer; no solver source was repaired or changed for these experiments.

\section{Reproduction and implications for implementation}
The accompanying directory contains the benchmark script, all seven input
sets, every raw repeat, first-call durations, counters, profiles, hashes,
derived ratios, and the magnitude checks. Sources are pinned to the commits
in the bibliography; byte hashes and hashes with normalized line endings are
recorded in \texttt{provenance.json}. The measured new sources equal the
packaged sources apart from line endings.

From the FastHVChan root, after installing the requirements in
\texttt{NumericalChanHVND} and checking out the pinned old repository at
\texttt{../HV4DMagnitude}, run the following single command (line breaks
are shown for readability):
\begin{quote}\small\ttfamily
python NumericalChanHV4D/benchmarks/prefix4/compare\_prefix4.py\\
\hspace*{1em}--old ../HV4DMagnitude --new NumericalChanHVND\\
\hspace*{1em}--reference . --out prefix4-reproduction --profiles
\end{quote}
The separate \texttt{check\_prefix4\_magnitude.py} accepts
\texttt{--old ../HV4DMagnitude --out magnitude-checks.json}.

The old compiled $2+2$ evaluator is the appropriate practical 4D baseline.
A faster new backend should keep a whole contraction loop inside compiled code
and control term-state growth. The compression ablation motivates an adaptive
policy, but adopting one while retaining a proved worst-case bound requires
additional analysis. These are concrete follow-up targets, not improvements
implemented or new complexity theorems established by this report.

\begin{thebibliography}{9}
\bibitem{old} M. Emmerich, \emph{HV4DMagnitude}, source snapshot
\texttt{0c668e7}, inspected 4 October 2026.
\href{https://github.com/emmerichmtm/HV4DMagnitude/tree/0c668e76520201f41ac0f144caa8d468d6c84e86}{Pinned repository}.
Relevant files: \texttt{solver.py}, \texttt{prefix4.py}, \texttt{fastkernel.py},
and \texttt{notes/audit.md}.
\bibitem{chan} M. Emmerich, \emph{FastHVChan}, reference source snapshot
\texttt{08539aa}, inspected 4 October 2026.
\href{https://github.com/emmerichmtm/FastHVChan/tree/08539aa99e94e181d7230a8288885b3294621c9e}{Pinned repository}.
The measured Python files are \texttt{chan\_orthant\_dby3.py} and
\texttt{chan\_hypervolume.py}. These are repository implementations of the
Chan approaches, not a claim to benchmark code authored by Chan.
\bibitem{new} \emph{NumericalChanHV4D and NumericalChanHVND},
FastHVChan commit \texttt{97012f7}.
\href{https://github.com/emmerichmtm/FastHVChan/tree/97012f7/NumericalChanHVND}{Pinned shared numerical implementation}.
The 4D Numba wrapper selects dimension four in that shared engine.
\end{thebibliography}
\end{document}
